Consider now the $\tau^{-1}$-linear operator

\begin{equation}\tag{4.5.13}
\tau^{-1}\circ\psi_p\circ H(X):L(b)\longrightarrow L(b).
\end{equation}

Its iterates are

\begin{equation}\tag{4.5.14}
(\tau^{-1}\circ\psi_p\circ H(X))^a
=\tau^{-a}\circ\psi_{p^a}\circ
H(X)H^\tau(X^p)\cdots H^{\tau^{a-1}}(X^{p^{a-1}}),
\end{equation}

and, if $\operatorname{ord}(q)$ divides $a$, then $\tau^{-a}=\operatorname{id}$.
Define the $\tau^{-1}$-linear operation

\begin{equation}\tag{4.5.15}
\left\{
\begin{aligned}
\alpha:L^0\!\left(\frac1{p-1}\right)
&\longrightarrow L^0\!\left(\frac1{p-1}\right),\\
\alpha&=\frac1p\,\tau^{-1}\circ\psi_p\circ H(X).
\end{aligned}
\right.
\end{equation}

Then (4.5.12) takes the form

\begin{equation}\tag{4.5.16}
\operatorname{card}V^*(\mathbb F_{q^\nu})
=(q^\nu-1)^{n+1}
+(q^\nu-1)^{n+2}
\operatorname{tr}\!\left(
\left.\alpha^{\operatorname{ord}(q^\nu)}
\right|L^0\!\left(\frac1{p-1}\right)\right).
\end{equation}

Put

\begin{equation}\tag{4.5.17}
S=\{1,\ldots,n+2\}.
\end{equation}

For every nonempty subset $A\subset S$, define the subspace

\begin{equation}\tag{4.5.18}
\begin{aligned}
L^A(b)=\{\,&\text{series $\sum_U B_UX^U\in L(b)$ such that, for every
monomial $X^U$ with $B_U\ne0$,}\\
&\text{the indices $i$ for which $U_i>0$ are exactly $0$ and the
$j\in A$}\,\}.
\end{aligned}
\end{equation}

Then

\begin{equation}\tag{4.5.19}
L^0(b)=\bigoplus_{\substack{A\subset S\\A\ne\varnothing}}L^A(b).
\end{equation}

\medskip
\noindent\emph{Source note.}
The print drops $\nu$ from the factor $q^\nu-1$ in (4.5.10) and from the
denominator $q^\nu$ in (4.5.12), although both formulas are direct rewrites
of (4.4.6). It also transposes the formula number (4.5.14) as (5.4.14).
The diplomatic French preserves all three printed forms; corrected French
and English use the algebraically forced readings displayed above.
